\section{Related Work}
\label{sec:related}

% ---------------------------------------------------------------- method summary
\begin{table}[t]
\centering
\small
\setlength{\tabcolsep}{3pt}
\begin{tabular}{@{}llll@{}}
\toprule
Method & Threshold $q_y$ from & If $n_y=0$ & Specific \\
\midrule
\textsc{Standard} \cite{vovk2005algorithmic}
  & one shared quantile      & shared quantile          & n/a \\
\textsc{Classwise} \cite{sadinle2019least}
  & own-class quantile       & undefined, $q_y=\infty$  & no \\
\textsc{Clustered} \cite{ding2023class}
  & cluster quantile         & null cluster, shared     & no \\
\textsc{Rc3p} \cite{shi2024rank}
  & quantile + rank budget   & undefined, $q_y=\infty$  & no \\
\textsc{Fuzzy} \cite{ding2026conformal}
  & kernel-weighted quantile & collapsed, shared        & no \\
\textsc{Pas} \cite{ding2026conformal}
  & one shared quantile      & shared quantile          & n/a \\
\textsc{Interp-Q} \cite{ding2026conformal}
  & $\tau q_y^{\mathrm{cw}}+(1-\tau)\hat q$ & capped, then blended & weakly \\
\textsc{Tacp} \cite{liu2026tail}
  & shared quantile, head penalty & head/tail only      & by group \\
\midrule
\textbf{PCC} (ours)
  & $g_\theta(\varphi(y))$ from weights & predicted     & \textbf{yes} \\
\bottomrule
\end{tabular}
\caption{Methods that give a class its own conformal threshold. All of them target at least
marginal coverage; the columns that separate them are what happens to a class with no labelled
calibration examples, and how specific to that class the threshold it receives can be. Three
methods do produce a defined threshold there without per-class labels: \textsc{Pas} reshapes the
score using prevalence, \textsc{Interp-Q} caps the undefined classwise quantile at one before
blending, and \textsc{Tacp} penalises the score of head classes, which gives one degree of
freedom per prevalence group instead of one per class. Behaviours were obtained by running the
authors' released implementations on the same held-out classes, except for \textsc{Tacp}, whose
authors release no code; \cref{tab:competitors} reports the coverage that results.}
\label{tab:methods}
\end{table}

\paragraph{Class-conditional conformal prediction.} Classwise conformal prediction, an
instance of Mondrian conformal prediction \cite{vovk2005algorithmic}, calibrates an
independent quantile per class and thereby attains class-conditional coverage, at the price of
very large sets when classes are thin \cite{sadinle2019least,lofstrom2015bias}. Two lines of
work soften that price by sharing information between classes: \textsc{Clustered} CP groups
classes by an embedding of their score distributions and calibrates one quantile per cluster
\cite{ding2023class}, while rank-calibrated CP couples a per-class quantile with a per-class
rank budget \cite{shi2024rank}. Most recently, \citet{ding2026conformal} embed each class by
the quantiles of its own score distribution and compute a kernel-weighted quantile that
borrows from nearby classes, alongside a prevalence-adjusted score function.

That last method is the closest to ours in structure, and the difference between them is what
this paper is about. Their embedding of class $y$ is computed from the labelled calibration
scores of class $y$, so when a class has none the embedding is undefined and their
implementation assigns it the extreme value of the embedding range, with the consequence that
every zero-example class receives one shared threshold. Our descriptors are computed from the
classifier's weights and never touch a label, so a class with zero examples still receives a
threshold of its own. \Cref{tab:competitors} reports both behaviours measured on the same
data.

\paragraph{Adaptive conformal prediction and score design.} Normalised or locally adaptive
conformal prediction predicts a per-instance notion of difficulty from the features and
rescales the score by it \cite{lei2014distribution}, which is the per-instance counterpart of
what we do per class. A parallel line designs the score function itself so as to trade set
size against conditional coverage, including APS \cite{romano2020classification}, RAPS
\cite{angelopoulos2021uncertainty} and SAPS \cite{huang2024conformal}; we report results for
all three in \cref{tab:depth}. What is new in our setting is not the idea of modulating a
conformal quantity by side information but the side information itself, which remains
available for classes the calibration set never saw.

\paragraph{Prevalence as the source of the correction.} A second family leaves the per-class
quantile alone and reshapes the score using class frequency. \citet{ding2026conformal} divide
the softmax probability by an estimate of the class prior, which is optimal for macro-coverage
among procedures of a given expected set size, and blend the classwise and marginal quantiles
with a user-chosen weight. \citet{liu2026tail} add a rank-dependent penalty to the scores of
head classes, with a smooth reweighting in their sTACP variant, and prove a smaller head-tail
coverage gap. Both are defined for a class with no calibration examples, which makes them the
closest competitors to ours that remain applicable in our regime, and both draw their per-class
signal from a single scalar. Prevalence turns out to be the weakest of the descriptors we
tried: it is significantly worse than geometry on its own (\cref{tab:ablation}) and is beaten
by a random projection of the same dimension (\cref{tab:competitors}), which is the empirical
reason we look at the classifier's geometry instead.

\paragraph{Group-conditional coverage.} A third family asks for coverage conditional on
membership of a group instead of treating classes one at a time. \citet{bairaktari2025kandinsky} give the most
general form of this, allowing overlapping and fractional group memberships defined jointly on
covariates and labels, obtained from a single quantile regression over the span of the group
functions. Our descriptors could serve as such group functions, which would replace our
empirical claim for zero-example classes with a group-conditional guarantee, at the price of a
guarantee about neighbourhoods in $\varphi$-space instead of about individual classes. We regard
this as the most promising route to a formal statement in the $n_y=0$ regime and do not pursue
it here.

\paragraph{Conformal prediction under distribution shift.} Coverage can be restored under
covariate shift by importance-weighting the calibration scores \cite{tibshirani2019conformal},
or over time by adapting the miscoverage level from observed errors \cite{gibbs2021adaptive}.
Both require something our setting does not provide, namely a density ratio or online label
feedback, so we cite them as the route to a guarantee and read our corrupted-input experiment
as a stress test, making no coverage claim there (\cref{app:imagenetc}).

\paragraph{Long-tailed recognition.} Reweighting and logit adjustment improve accuracy under
class imbalance \cite{menon2021long}, and fine-grained biodiversity datasets have made the
regime concrete \cite{garcin2021plantnet,vanhorn2018inaturalist}. Our concern is complementary
and post-hoc: given whatever classifier such methods produce, how should its prediction sets
be calibrated for the classes that training consumed.
